\DocumentMetadata{
  pdfversion=2.0,
  pdfstandard=ua-2,
  lang=en-US,
  tagging=on,
  tagging-setup={math/setup=mathml-SE,tabsorder=structure}
}
\documentclass[11pt,letterpaper]{article}
\usepackage[margin=0.68in,includefoot]{geometry}
\usepackage{newcomputermodern}
\usepackage{amsmath,amscd,mathtools}
\providecommand{\square}{\mdlgwhtsquare}
\usepackage[hidelinks]{hyperref}
\hypersetup{
  pdftitle={Complex Analysis PhD Exam, June 1997},
  pdfauthor={University of Florida Department of Mathematics},
  pdfsubject={Graduate mathematics examination},
  pdfdisplaydoctitle=true,
  bookmarksopen=true
}
\setcounter{secnumdepth}{0}
\setlength{\parindent}{0pt}
\setlength{\parskip}{0.28em}
\setlength{\emergencystretch}{3em}
\setlength{\leftmargini}{2.15em}
\setlength{\leftmarginii}{2.65em}
\setlength{\leftmarginiii}{3em}
\renewcommand{\labelenumi}{\textbf{\theenumi.}}
\pagestyle{empty}
\raggedbottom
\allowdisplaybreaks
\newenvironment{examproblems}[1]{%
  \begin{enumerate}%
  \setcounter{enumi}{#1}%
  \setlength{\topsep}{0.35em}%
  \setlength{\partopsep}{0pt}%
  \setlength{\itemsep}{0.48em}%
  \setlength{\parsep}{0.18em}%
}{\end{enumerate}}
\newenvironment{examsubparts}{%
  \begin{enumerate}%
  \setlength{\topsep}{0.18em}%
  \setlength{\partopsep}{0pt}%
  \setlength{\itemsep}{0.16em}%
  \setlength{\parsep}{0.1em}%
}{\end{enumerate}}
\newenvironment{examinstructions}{%
  \par\begingroup\itshape
}{%
  \par\endgroup\vspace{0.18em}
}
\makeatletter
\renewcommand\section{\@startsection{section}{1}{0pt}%
  {-1.25ex plus -.25ex minus -.1ex}%
  {0.55ex plus .1ex}%
  {\normalfont\large\bfseries}}
\renewcommand{\@maketitle}{%
  \newpage\null\vskip -1.2em%
  {\footnotesize\bfseries
    \href{https://www.ufl.edu/}{University of Florida}, \href{https://math.ufl.edu/}{Department of Mathematics}\par}%
  \vskip 0.18em%
  \tagstructbegin{tag=Title}%
    {\LARGE\bfseries\href{https://gma.math.ufl.edu/exams/complex-analysis-phd/}{Complex Analysis PhD Exam}, June 1997\par}%
  \tagstructend%
  \vskip 0.35em%
  \tagmcbegin{artifact}%
    \hrule height 1pt%
  \tagmcend%
  \vskip 0.55em%
}
\makeatother
\title{Complex Analysis PhD Exam, June 1997}
\author{}
\date{}
\begin{document}
\maketitle
\thispagestyle{empty}
\begin{examinstructions}
Do each of the problems.
\end{examinstructions}
\begin{examproblems}{0}
\item
Evaluate the integral \(\displaystyle \int_0^\infty \frac{x^{a-1}}{1+x}\,dx\), \(0<a<1\).
\item
Suppose \(\Omega=\{z:-1<\operatorname{Re}z<1\}\). Find an explicit formula for the 1-1 conformal mapping~\(f\) of~\(\Omega\) onto \(\{z:|z|<1\}\) for which \(f(0)=0\) and \(f'(0)>0\). Compute \(f'(0)\).
\item
Let~\(F\) be the class of all analytic functions~\(f\) in \(\mathbb D=\{z:|z|<1\}\) for which 
\[
\iint_{\mathbb D}|f(z)|^2\,dx\,dy\leq 1.
\]
 Is this a normal family? Justify your answer.
\item
Let~\(g\) be a polynomial. Show that 
\[
2\iint_{\mathbb D}(1-|z|^2)|g'(z)|^2\,dx\,dy\leq \int_0^{2\pi}|g(e^{i\theta})|^2\,d\theta,
\]
 where \(\mathbb D=\{z:|z|<1\}\).
\item
This problem has three parts.
\begin{examsubparts}
\item[\textnormal{(a)}]
Give a definition of subharmonic functions.
\item[\textnormal{(b)}]
If~\(f\) is analytic in~\(\Omega\), prove that \(\ln(1+|f(z)|^2)\) is subharmonic.
\item[\textnormal{(c)}]
Let~\(f\) be analytic in~\(\Omega\). Is \(\ln|f(z)|\) subharmonic in~\(\Omega\)? Justify your answer.
\end{examsubparts}
\item
Construct an entire function whose zero set is the set of Gaussian integers \(a+bi\), where \(a,b\in\mathbb Z\).
\item
For \(|z|<1\), define \(f(z)\) by 
\[
f(z)=z+\frac{z^3}{3}+\frac{z^5}{5}+\frac{z^7}{7}+\cdots.
\]
 Extend \(f(z)\) analytically in a suitable way so that it can be defined at \(z=i\). Using your analytic continuation, determine the radius of convergence of the Taylor series of~\(f\) with \(z=i\) as center.
\item
Show \(ze^z=1\) has infinitely many solutions.
\item
Define the order and genus of an entire function. What is the relation between the order and genus? Find the order of \(\sin z\).
\end{examproblems}
\end{document}
